\chapter{The example driver}\label{ch:driver}

\texttt{examples/driver.py} runs a representative set of the public API on small synthetic
inputs and prints the results. \texttt{examples/driver\_expected.txt} holds its output on a CPU,
and \texttt{tests/test\_example\_driver.py} compares the two. The driver is the quickest way to
see that an installation works and to compare it with the expected values.

\section{Running it}\label{driver:running}

From the root of a source checkout, with Exacta installed:
\begin{lstlisting}[language=bash]
python examples/driver.py                              # CPU
python examples/driver.py --device cuda                # the same on a CUDA device
python examples/driver.py 2> /dev/null | diff - examples/driver_expected.txt
\end{lstlisting}
It takes a few seconds, most of it spent importing PyTorch. The standard output is the result;
a line with the versions of Exacta, PyTorch, NumPy and Python, the platform, the device and the
number of threads goes to standard error, so that the standard output does not depend on the
machine. The exit status is zero when every check at the end of the output passes.

\section{What it does}\label{driver:what}

The driver draws no random numbers. Every input is built from integer arithmetic: two cases of
$22\times26\times24$ voxels whose target is a discrete ball and whose predictions are the ball
moved by three voxels along the last axis (plus a small separate ball), and a larger ball. The
spacing is $(1.2, 0.4, 0.4)$ millimetres: anisotropic, and none of its values is exactly
representable in binary, so products such as $3\times0.4$ are inexact in binary64. The sections are:

\begin{enumerate}[nosep]
\item the inputs;
\item the exact distance transform of \texttt{exacta.losses.distance\_maps} on a
  $5\times6\times7$ grid with spacing $(1.3,0.3,0.7)$, compared bit for bit with a brute-force
  search written in the driver;
\item HD, HD95 and ASSD with the three surface models (\texttt{voxels}, \texttt{surfels},
  \texttt{mesh}), and the overlap scores and counts;
\item NSD at a tolerance of 1.2 mm, where distances tie with the tolerance, with the default tie
  margin and with \texttt{tie\_rtol=0};
\item the presets that reproduce MONAI, MedPy, DeepMind's \texttt{surface-distance} and
  MeshMetrics (Chapter~\ref{ch:tools});
\item a stateful metric (\texttt{HD95}) and its \texttt{per\_case}, \texttt{per\_class} and
  \texttt{mean}, and the reporting signature of a score, parsed back and recomputed to the
  same bits;
\item \texttt{DiceLoss} and \texttt{DiceCELoss} with value and gradient norm, and the check that
  the second is the first plus PyTorch's cross-entropy.
\end{enumerate}

Numbers are written with Python's \texttt{repr}, which prints a float64 with the digits needed to
read it back bit for bit. Lines without a marker come from computations that Exacta rounds or
counts exactly; they are the same on every platform, thread count and device. Lines that end
with \texttt{(approx)} are the losses, which run on PyTorch operations or fused Triton kernels, and
ASSD and NSD with mesh weights, which sum triangle areas that PyTorch computes (a cross product
and a norm). The last bits of both depend on the PyTorch build and the device: the mesh ASSD
differed by one unit in the last place on Linux aarch64 with PyTorch 2.15, and not with 2.4.1.
The test compares these lines to a relative tolerance of $10^{-9}$ and everything else as text.

\section{Reading the output}\label{driver:reading}

The output below was produced on x86-64 Linux with PyTorch 2.14.1 on the CPU.
Three lines show what the conventions of Chapter~\ref{ch:conventions} do:
\begin{itemize}[nosep]
\item \texttt{3 * 0.4} is \texttt{1.2000000000000002}, so the 1.2 mm distance of three voxels
  lies above a tolerance of 1.2 in binary64.
\item The voxel NSD of case 0 is \texttt{0.9717} with the default margin and \texttt{0.7778} with
  \texttt{tie\_rtol=0}: the surface elements that lie exactly on the tolerance are counted in
  the first and left out in the second. The surfel and mesh values move in the same way.
\item HD95 of case 0 is \texttt{3.2242} with \texttt{combine="max"} (MONAI), but with
  \texttt{combine="pooled"} (MedPy) it is \texttt{1.2000000000000002}: one name, two quantities.
\end{itemize}

\lstinputlisting[language=output,basicstyle=\ttfamily\scriptsize,caption={Output of
\texttt{python examples/driver.py} on the CPU (\texttt{examples/driver\_expected.txt}).}]{../../examples/driver_expected.txt}

\section{How it was checked}\label{driver:checked}

The output was identical in two runs, and with \texttt{OMP\_NUM\_THREADS} set to 1 and to 8,
including the loss lines. On a CUDA device (\texttt{--device cuda}) every line without the
marker was identical to the file, and the two gradient norms differed from it in the last
digit. The test file checks the same properties: the output equals the expected file, it does
not depend on the thread count, and, when a GPU is present, the exact lines are the same on it.
After a deliberate change of the driver, regenerate the expected file with
\begin{lstlisting}[language=bash]
python examples/driver.py 2> /dev/null > examples/driver_expected.txt
\end{lstlisting}
and review the difference.
